\scanpage{107}
\uline{Automorphismes des \deleted{complexes} \inserted{graphes} trialitaires}.

Soit $(S,A)$ un \deleted{complexe} \inserted{graphe} trialitaire, $\underline t=\{a,b,c\}$ un triangle. On a défini une bijection involution \deleted{$\unsure{\sigma_t|\sigma_{\underline t}}$ de $\unsure{S|S,S}$}

\[ \sigma_{\underline t}=\sigma_{\{a,b,c\}}: S\mathrel{\unsure{-|\to}}S \]

par

\[
\begin{aligned}
\sigma_{\underline t}\mid\underline t&=\mathrm{id}_{\underline t}\\
\sigma_{\underline t}\mid E_a^{\unsure{t|\underline t}}\deleted{\unsure{(t)}}&=\sigma_a^{\underline t}
\qquad\text{(transforme $x\in E_a$ en l'unique}\\[-2pt]
&\hspace{7em}\text{él de $E_a$ lié à $x$)}
\end{aligned}
\]

\uline{Proposition} \qquad a) $\sigma_{\underline t}$ est un automorphisme du \deleted{complexe} graphe $(S,A)$, et $\underline t$ est l'ens des pts fixes de $\sigma_{\underline t}$.

b) L'application $\underline t\longmapsto\sigma_{\underline t}$ de $T$ dans $G=\mathrm{Aut}(S,A)$ est injective, et a comme image l'ens des autom involutifs \inserted{$u$} de $(S,A)$ tels que \deleted{l'ens des} $\operatorname{card} S^u=3$

C'est $\pm$ évident\ldots

Soit $G\unsure{=|\neq}\mathrm{Aut}(S,A)$. Les $\sigma_t$ sont $\in G$. Soit $G^{\unsure{0|\circ}}$ le grpe qu'ils engendrent. Si $g\in G$, on a

\[ \operatorname{int}(g).\sigma_t=g\,\sigma_t\,g^{-1}=\sigma_{g(t)} \]

donc $G^0$ est un ss. grpe invariant.

% archive sheet number: 106
